\scanpage{73}
\deleted{\unsure{pour une section|pour ces sections}} considérons (pour $S,X,t$ fixés) la catégorie des $(X',t',\underline G)$, \deleted{où $X'$ est un $X$-schéma} \inserted{tels que dessus, avec} comme morphismes
\[
\unsure{u:|u<}\ (X',t',G)\longrightarrow(X'_1,t'_1,G_1)
\]
% slanted
\margin{\uline{NB} On n'exige pas que $n=n_1$. Pour $X',t',G$ fixés, $(X'_1,t'_1,G_1)$ est déterminé à isom can. près, quand on se donne $G\to G_1$}
les \deleted{\unsure{hom.|homs}} couples \inserted{de} $X'\to X'_1$ ($X$-morphismes et de $G\to G_1$, compatibles dans un sens évident avec les sections marquées et les opérations. Considérons d'autre part les \deleted{triples} systèmes
\[
(\underline G,\deleted{\unsure{\xi}}n,\deleted{\unsure{\xi}}i,\deleted{\unsure{n|i}}\xi,\deleted{\unsure{k}}T,k)
\]
où $n\in\mathbb N^\ast$, \inserted{$\underline G$ un schéma en groupes fini étale sur $S$,} $\xi\in H^0(S,R^1f^\circ_\ast f^{\circ\ast}(\underline G))$ tels que les rev. locaux (\unsure{ét. loc.|étale.|étales} sur $S$) de $X$ qu'elle décrit soient modérément ramifiés d'ordre $n$ le long de $t$, $i:\mu_{n\,S}\hookrightarrow\underline G$ est un mono. de groupes, $T$ un Module inv. sur $S$, et $k$ un isom
\[
T^{\otimes n}\xrightarrow{\sim} T_{X/S,t}
\]
avec la définition ad hoc des morphismes (\uline{NB} \unsure{bien|bête|bien entendu} un $\underline G\to\underline G_1$ doit induire un \unsure{épimorphisme|homomorphisme|\deleted{épi}morphisme} $\mu_{n\,S}\to\mu_{n_1\,S}$, qui soit l'hom \deleted{\unsure{gén.|groupe|qui}} $\lambda\longmapsto\lambda^{n/n_1}$ \ldots). On a un foncteur canonique de la première catégorie dans la
% archive sheet number: 72
